\section{The least ordinary-reduction field in a tensor case}
\label{sec:original-a1}

\begin{theorem}\label{thm:original-a1}
Let \(A/K\) be a polarized abelian variety over a number field. Suppose
\(G^{\mathrm{ad}}\), for its Betti Mumford--Tate group \(G\), is
\(\Q\)-simple and has geometric type \(A_1\) in every factor. Suppose a
Rosati-compatible CM field \(E\) acts over \(K\), and each geometric
\(E\)-block of \(H^1\) is irreducible under \(G^{\mathrm{der}}\).
Let \(R\) be the centroid field of the derived Lie algebra. Write
\(S\subset\operatorname{Hom}_{\Q}(R,\overline{\Q})\) for the embeddings
where the Hodge cocharacter projects nontrivially. Assume \(|S|\le2\).
Then the stabilizer of \(S\) under postcomposition in \(G_K\) is the
kernel of its action on the centroid. It defines a finite Galois
extension \(K_S/K\). Outside a density-zero set, the
ordinary places of \(K\) are precisely those that split completely in
\(K_S\). Thus their natural density is \([K_S:K]^{-1}\), and \(K_S\)
is the least extension over which the ordinary density is one.
\end{theorem}

The one-active tensor case has the familiar binomial Newton polygons
of Mumford type \cite[\S2]{nootrealization}\cite[Theorem~1.2]{shengzuo}.
With two active labels, one adjoint block can lose its gap while the
other retains it. The trace of a tensor block recovers the product of
the projective traces of its supporting factors, so the two local
valuations can be compared before taking a power of Frobenius.

\begin{proof}
Fix the Betti realization \(V=H^1\), and write \(C_\lambda\) for the
set of simple factors supporting the block of an embedding \(\lambda\)
of \(E\). Over an algebraic closure an irreducible block lifts to a
tensor of symmetric powers of the two-dimensional standards of its
support. A block meets at most one active label and has standard
exponent there. Indeed each nontrivial active exponent \(a\) adds
\(a\) to the difference between the number of distinct Hodge weights
and one, as a lower bound obtained by adding the ordered weight sets.
The only Hodge weights on \(V\) are zero and one. Galois transitivity
on the embeddings of \(R\) and equivariance of supports and exponents
then make every factor in every support standard. All supports have
the same size \(k\), because the block dimensions are equal. Write
\(h=2^k\) for this dimension. Each label of \(R\) occurs in the same
positive number \(c\) of supports.

The trace pairing on \(\operatorname{End}(V)\) is nondegenerate on the
derived Lie algebra \(\mathfrak d\), since its simple factors are
orthogonal and each has nonzero trace index. The orthogonal projector
onto \(\mathfrak d\) and its centroid action are absolute Hodge
tensors \cite[Proposition~2.9 and Theorem~2.11]{delignehodge}.
Vasiu's Mumford--Tate comparison in type \(A_1\)
\cite[Theorem~1.3.4(a) and Example~1.3.5(iv)]{vasiu}
shows that full auxiliary monodromy normalizes \(\mathfrak d\), so it
fixes the projector and acts on the centroid by field automorphisms
\(\beta_g\). This is the same action on the absolute Hodge tensors in
each realization. If \(\theta_g\) denotes postcomposition by \(g\) on
the embeddings of \(R\), the semilinear action on de Rham cohomology
sends the \(\rho\)-eigenspace to that for \(\theta_g(\rho)\circ
\beta_g^{-1}\). It preserves the positive adjoint Hodge filtration
and hence \(S\). Since precomposition commutes with postcomposition,
\(\theta_g\) preserves \(S\) exactly when precomposition by
\(\beta_g\) does. Because \(E\) acts over \(K\), the linear Galois
action fixes every coefficient block. Its permutation of derived
factors preserves each support \(C_\lambda\). If it preserves \(S\),
it fixes each active label, since every such label occurs in a
support with no other active label. A field automorphism fixing
one embedding is the identity. This identifies the stabilizer
with \(\ker\beta\).

Take a good absolute degree-one place \(v\) with arithmetic Frobenius \(g\).
Omit the finite bad set for the coefficient and tensor fields.
Pass to a fixed finite extension defining the tensors and level and
containing the reflex field. The local extension is unramified at
\(v\). In the integral model of Section~\ref{sec:adjoint}, the lift
gives a geometric special point \(s\) whose underlying abelian variety
is the base change of \(A_v\). Proposition~1.3.6 and \S1.3.7 of
\cite{kmps} construct the crystalline tensors at \(s\), fixed by the
one-step Witt-semilinear Frobenius \(\Phi\). In particular, their
tensor trivialization gives \(\Phi=b\sigma\) with \([b]\in B(G_{\Q_p})\),
even when \(s\) is defined over a larger residue field. The
tensor comparison identifies the Hodge class of the lift with
the specified one \cite[Theorem~(0.3) and \S5]{blasius}.
Thus Lemma~1.3.8 of \cite{kmps} gives
\(\kappa_G(b)=\mu^\natural\) and the group Newton inequality
for this one-step operator. Since \(\Phi\) fixes each centroid tensor
and acts on its eigenvalues by \(\sigma\), it permutes the adjoint
labels by \(\theta_g\). The dominant adjoint Hodge gaps are one on
\(S\) and zero elsewhere. Their average on an orbit of length \(a\)
meeting \(S\) in \(m\) labels bounds the nonnegative Newton gap \(d\),
so \(ad\le m\).

Write \(d_\rho\) for the gap on label \(\rho\). The slopes on
\(V_\lambda\) are its Hodge mean shifted by all signed sums
\(\frac12\sum_{\rho\in C_\lambda}\epsilon_\rho d_\rho\), with
\(\epsilon_\rho\in\{-1,1\}\). The signs come from the standard tensor
weights. The central shift is the Hodge mean by the Kottwitz equality
on the abelianization, because the block Hodge multiplicities are
preserved by \(\theta_g\) when \(E\) acts over \(K\). If \(\theta_g\)
fails to preserve \(S\), take a block supporting an active label on a
partially active orbit. Its Hodge mean is \(1/2\). Ordinarity on that
block would require exactly one positive gap, equal to one, since
two positive gaps yield at least three distinct slopes. A gap of one
requires a fully active orbit. This block has just one Hodge-active
label and that label lies on a partially active orbit, so the
required gap cannot occur. Conversely, when \(\theta_g\) preserves
\(S\), attaining gap one at every active label gives ordinarity by
the same weight formula.

Now let \(\theta_g\) preserve \(S\), so \(\beta_g=1\). The adjoint
projector and centroid descend to the crystalline fiber over
\(\Q_p\), since inertia and Frobenius act trivially on them. Linear
crystalline Frobenius preserves this Lie algebra and acts linearly
over \(R\otimes\Q_p\). On extending scalars, semilinear Frobenius
is this linear map followed by scalar Witt Frobenius.

At an active orbit the local adjoint Lie algebra over its unramified
coefficient field is split by its positive Hodge root line, as in
Section~\ref{sec:adjoint}. Let \(B\) lift the local linear projective
action. Distinct eigenvalue valuations of \(B\) are integers by
the quadratic Newton polygon. Hence the gap at every active orbit
is zero or one, since \(d\le1\). For an active singleton, the
coefficient field is \(\Q_p\) and \(B\) can be chosen in
\(\mathrm{GL}_2(\Q_p)\). In this basis the semilinear projective
Frobenius is \(B\sigma\), so Kottwitz gives its determinant valuation
odd. The local calculation of Section~\ref{sec:adjoint} gives
valuation at least one at a flat active singleton and minus one
at a nonflat singleton for \(T=(\operatorname{Tr}B)^2/\det B\).
At any flat label, the invariant \(T=1+\operatorname{Tr}(\operatorname{Ad} B)\)
has nonnegative valuation, as the adjoint eigenvalue valuations are zero.

In both realizations linear Frobenius normalizes the derived group and
fixes its factors individually. Their automorphisms are inner and
Frobenius preserves the \(E\)-blocks. On a block its action is
therefore a scalar times \(\bigotimes_{\rho\in C_\lambda}B_\rho\), by
Schur's lemma. If \(a_\lambda,\delta_\lambda\) are the block trace and
determinant, this gives
\[
 U=\prod_\lambda\frac{a_\lambda^h}{\delta_\lambda}
  =\prod_\rho T_\rho^{ch/2}.
\]
The relative polynomial comparison in Section~\ref{sec:coefficients}
makes \(U\) rational with denominator a power of \(p\), and the Weil
bound makes it bounded independently of \(v\). If all active labels
are flat, it is integral at \(p\). Otherwise a failure can only be
one flat active singleton and one nonflat active singleton, since
there are at most two active labels. Their valuations also make
\(U\) integral at \(p\).
On every component inducing the identity on the centroid, the regular
function \(U\) is nonconstant. Multiply a representative by a derived
element to centralize the derived group. The block scalars then cancel
in \(U\), and varying one simple coroot varies just its \(T_\rho\).
Algebraic Chebotarev rules out the finitely many integer levels. This
proves the claimed ordinary criterion. It applies after every finite
extension of \(K\) with the same label set and the restricted Galois
action, so Chebotarev also yields the least extension assertion.
\end{proof}
